\begin{table*}[t]
\centering\footnotesize
\caption{The decided rules and their verdicts. Each rule was fixed with its threshold before any
result of the study existed (Section~\ref{sec:study}); a rule that fires removes or restricts the
claim it guards, and one that does not fire leaves it standing. The literature rule F4 fired
before the data and narrowed the claim on compensation to claim~2$'$
(Section~\ref{sec:remedies:comp}). F2 fires if neither primary cell shows a rise in $U$
over $\theta = 0$ of at least 1 percentage point (pp) with a 97.5\,\% interval above 0, and
restricts the diagnosis to the cell that shows it; F3 fires if 0.90 lies outside the Wilson
interval where the rise is shown; F5a and F5b if the mean $U$ exceeds 2\,\%; F6 if the mean error
of the predicted width ratio lies outside $\pm 10$\,\%; F7 if a baseline is valid and either it is
narrower or the loop is not valid; F8 if the 95\,\% interval of the paired difference in IACT does not lie wholly
below 0, and its step-size statement needs a larger compensated step in at least two operating
points and a smaller one in none. F1 is secondary and decides nothing.}
\label{tab:rules}
\input{\figdir/n1_table_rules.tex}
\end{table*}

\section{Pre-registered study}\label{sec:study}

The design below, the rules of Table~\ref{tab:rules} with their thresholds, the analysis plan and
the baseline configurations were frozen in a public repository on 29 October 2026, before any
result of the study existed (Data and code availability). Stage A of the run read training and
calibration rows only: it fitted the point models and compensators, tuned every step size, configured the
baselines and recorded every predicted width ratio, and its records were committed and pushed
before stage B read an evaluation row. Stage B recomputed stage A's quantities from the same files
and passed all 1359 checks against them exactly before evaluating. One deviation was logged before
any result: the arm that updates once per analyzer value seeds its window with 40 labels instead
of 200, because its calibration hours hold only 160; no rule is decided on that arm. Claude
(Anthropic) assisted in writing the analysis code and in screening the literature; the code's
results are checked as above, and every screening verdict is recorded with its source in the
repository.

\subsection{Data and point models}\label{sec:study:data}

\emph{Tennessee Eastman process (TEP).} The plant of \citet{downs1993} was simulated with the closed-loop
Fortran code of Russell, Chiang and Braatz \citep{russell2000}, with the feed-composition,
feed-temperature and reaction-kinetics disturbances IDV(8), IDV(10) and IDV(13) active
throughout, at three operating points set by the reactor-feed composition loops: the base case of
\citet{downs1993} and two with more G in the product (G at about 53.8, 58.2 and 63.4~mol\,\%;
operating points 1 to 3, labelled mode 1 to 3 in the figures). Each run drew its random
numbers from its own MRG32k3a stream \citep{lecuyer1999}, because the simulator's generator loses
precision in double-precision arithmetic \citep{kitchin2026} and then cycles within 0.55 to
141~h (supplementary material). For each operating point, the first five streams in a fixed
order whose 100-h run did not shut down were kept (5 of 26 candidates survived at the first, 14 of
25 and 17 of 25 at the others), so the evidence is conditional on runs that stay within the
plant's limits. Rows are 60~s apart: hours 0--30 train the point model, 30--70 calibrate and
70--100 evaluate. The target is the product G analysis, XMEAS(40). The point model is a ridge
regression on the 22 continuous measurements and 12 manipulated variables taken 15~min earlier,
its penalty chosen by fivefold cross-validation with the one-standard-error rule. One model per
run serves three arms that differ only in the calibrator: T1 updates every 3~min
($\theta = 0, 2, 4$), T2 every minute ($\theta = 0, 4, 14$) and T3 once per new analyzer value
($\theta = 0$). Gradient boosting on the same inputs checks the estimator class at T1.

\emph{Debutanizer.} The public data of \citet{fortuna2007} are one stream of 2394 rows. The butane
content is regressed on input $u_5$ 15 rows earlier by least squares on the first 1675 rows, whose
1660 residuals form the calibration segment; the other 719 rows are evaluated, at
$\theta = 0, 1, 3, 7$ and 12. The lag was chosen in earlier work on rows 0 to 2033, so for that one
choice the first 359 evaluated rows are not out of sample.

\begin{figure*}[t]
\centering
\includegraphics[width=\textwidth]{n1_diagnosis.pdf}
\caption{The diagnosis. Uninformative share $U$ (top) and coverage (bottom) against the dead time
$\theta$ for three ways of choosing the step: replay at the plant's $\theta$ to $\varepsilon/2$,
delay-blind tuning (replay at $\theta = 0$ to $\varepsilon$) and the white-noise rule
(Section~\ref{sec:analysis:persistence}); arms T1 (3-min updates; a, b) and T2 (1-min updates; c,
d), without and with compensation. Means over the 15 runs with 95\,\% $t$ intervals; markers are
offset around their $\theta$. Grey lines: the budget for $U$, 2\,\%, and the nominal coverage,
0.90. Coverage stays between 0.885 and 0.902 in every cell, while $U$ rises with $\theta$ under the two policies
that ignore the delay or the persistence.}
\label{fig:diagnosis}
\end{figure*}

\subsection{Cells, baselines and inference}\label{sec:study:cells}

A cell fixes the loop, the step-size policy, the pairing and the compensation; every contrast
varies one of them on a shared residual stream. The loops are the sequential loop of
Eq.~\eqref{eq:loop} (primary); split conformal, calibrated once (loop off); the level projected
onto $[0.01, 0.5]$; $\theta + 1$ interleaved delay-free chains \citep{elhalabi2026}; a
Smith-predictor form, in which the model $\mathrm{E}[e \mid \alpha] = \alpha$ swaps the miss rate
expected at the stored level for that at the current one; and decaying steps $c\,n^{-0.6}$
\citep{angelopoulos2024decay}. The policies are the fixed steps 0.005 \citep{gibbs2021aci} and
0.05, delay-blind tuning, replay to $\varepsilon$ and to $\varepsilon/2$, and the white-noise rule;
every loop is tuned by the same replay. The first 200 calibration scores seed the replay's window,
and every replayed update counts; on the Tennessee Eastman data the calibration segments of an
operating point's five runs are pooled, standing in for the longer history a plant would tune on.
Each run is evaluated from a window holding its last 200 calibration scores. Stored pairing is
primary; arrival pairing at the documented delay is the negative control.

Four baseline configurations get the same point model and only the labels in hand at issue.
Conformal PID control \citep{angelopoulos2023pid} runs with three scorecasters: quantile
regression on the three latest labelled scores, a linear form that \citet{areces2025} fit online;
the Theta model of the authors' code; and the compensator of Eq.~\eqref{eq:comp}, which gives it
the structure of AcMCP \citep{wang2024acmcp}. Its gains are tuned by the same replay to the same
target. SPCI \citep{xu2023spci} predicts the residual quantiles at horizon $\theta + 1$ with a
quantile random forest, without coverage feedback. Each follows its paper and, where the paper is
silent, the authors' code, with departures, each with its reason, listed in Appendix A of the
pre-registration (among them a warm start from the loop's seed rows and an unpenalised quantile
regression).

The 15 runs are independent replications: means carry 95\,\% $t$ intervals over the runs,
contrasts are paired by run, and pooled coverage carries a Wilson interval; the five runs of an
operating point share a tuned step, a dependence the $t$ intervals ignore. On the single
debutanizer stream, intervals come from a circular block bootstrap \citep{politis1992} of the
per-step indicators, with 2000 resamples, blocks of $\max(25, \lceil 2\,\mathrm{IACT}\rceil)$
steps and the same blocks for both arms of a contrast. The IACT is summed over the automatic
window of \citet{sokal1997} with $c = 5$. The rules are decided on the primary cells: F2 and F3 on
the uncompensated delay-blind loop, against $\theta = 0$ on the same stream, at T1, $\theta = 4$
and at the debutanizer's $\theta = 3$; F5 to F8 at T1, $\theta = 4$ (Table~\ref{tab:rules}).
Every other cell is descriptive. A power analysis on simulated AR(1) residuals set the margins: a
loop tuned to $\varepsilon$ met it on average, but its 15-run mean exceeded it in 21 to 38\,\% of
200 studies; tuned to $\varepsilon/2$, it exceeded the budget of $2\varepsilon$ in at most
4\,\%. F2 had a power of 0.91 to 0.99 in the Tennessee Eastman design, against 0.095 to 0.12 on
one stream of the debutanizer's length.

